
\documentclass{article}
\usepackage{colortbl}
\usepackage{makecell}
\usepackage{multirow}
\usepackage{supertabular}

\begin{document}

\newcounter{utterance}

\centering \large Interaction Transcript for game `matchit\_ascii', experiment `similar\_grid\_2', episode 4 with mistral{-}large{-}3{-}azure.
\vspace{24pt}

{ \footnotesize  \setcounter{utterance}{1}
\setlength{\tabcolsep}{0pt}
\begin{supertabular}{c@{$\;$}|p{.15\linewidth}@{}p{.15\linewidth}p{.15\linewidth}p{.15\linewidth}p{.15\linewidth}p{.15\linewidth}}
    \# & $\;$A & \multicolumn{4}{c}{Game Master} & $\;\:$B\\
    \hline

    \theutterance \stepcounter{utterance}  
    & & \multicolumn{4}{p{0.6\linewidth}}{
        \cellcolor[rgb]{0.9,0.9,0.9}{
            \makecell[{{p{\linewidth}}}]{
                \texttt{\tiny{[P1$\langle$GM]}}
                \texttt{U neemt deel aan een gezamenlijk raadspel.} \\
\texttt{Het doel is om erachter te komen of dit raster en een andere afbeelding die alleen ik kan zien, hetzelfde zijn.} \\
\texttt{Beschrijf eerst uw raster. Daarna geef ik mijn beschrijving en kunnen we elkaar vragen stellen over de rasters om erachter te komen of ze hetzelfde zijn. Begin nu uw korte rasterbeschrijving met "BESCHRIJVING:" gevolgd door de beschrijving. Voeg niets anders toe. Dit is uw raster:} \\
\texttt{X X X X X} \\
\texttt{X ▢ ▢ ▢ X} \\
\texttt{X ▢ ▢ ▢ X} \\
\texttt{X ▢ ▢ ▢ X} \\
\texttt{X X X X X} \\
            }
        }
    }
    & \\ \\

    \theutterance \stepcounter{utterance}  
    & \multicolumn{4}{p{0.6\linewidth}}{
        \cellcolor[rgb]{0.9,0.9,0.9}{
            \makecell[{{p{\linewidth}}}]{
                \texttt{\tiny{[P1$\rangle$GM]}}
                \texttt{BESCHRIJVING: Mijn raster is een 5x5 vierkant met een dikke rand van "X"{-}symbolen aan alle buitenste randen. In het midden bevindt zich een 3x3 gebied van lege vierkanten (▢), omringd door die rand.} \\
            }
        }
    }
    & & \\ \\

    \theutterance \stepcounter{utterance}  
    & & & \multicolumn{2}{p{0.3\linewidth}}{
        \cellcolor[rgb]{0.9,0.9,0.9}{
            \makecell[{{p{\linewidth}}}]{
                \texttt{\tiny{[GM$|$GM]}}
                \texttt{continue} \\
            }
        }
    }
    & & \\ \\

    \theutterance \stepcounter{utterance}  
    & & \multicolumn{4}{p{0.6\linewidth}}{
        \cellcolor[rgb]{0.9,0.9,0.9}{
            \makecell[{{p{\linewidth}}}]{
                \texttt{\tiny{[GM$\rangle$P2]}}
                \texttt{U neemt deel aan een gezamenlijk raadspel.} \\
\texttt{Het doel is om erachter te komen of dit raster en een andere afbeelding die alleen ik kan zien, hetzelfde zijn.} \\
\texttt{Beschrijf eerst uw raster. Daarna geef ik mijn beschrijving en kunnen we elkaar vragen stellen over de rasters om erachter te komen of ze hetzelfde zijn. Begin nu uw korte rasterbeschrijving met "BESCHRIJVING:" gevolgd door de beschrijving. Voeg niets anders toe. Dit is uw raster:} \\
\texttt{X ▢ X X X} \\
\texttt{X ▢ ▢ ▢ X} \\
\texttt{X ▢ ▢ ▢ X} \\
\texttt{X ▢ ▢ ▢ X} \\
\texttt{X ▢ X X X} \\
            }
        }
    }
    & \\ \\

    \theutterance \stepcounter{utterance}  
    & & & \multicolumn{4}{p{0.6\linewidth}}{
        \cellcolor[rgb]{0.9,0.9,0.9}{
            \makecell[{{p{\linewidth}}}]{
                \texttt{\tiny{[GM$\langle$P2]}}
                \texttt{BESCHRIJVING: Een 5x5 raster met een omlijning van X'en, behalve de eerste kolom die volledig uit X'en bestaat. In het midden van het raster (kolommen 2 t/m 4, rijen 2 t/m 4) zijn lege vakjes (▢), met uitzondering van de rechterbovenhoek (rij 1, kolom 4) die een X heeft. De rechteronderhoek (rij 5, kolom 4) is ook een X.} \\
            }
        }
    }
     \\ \\

    \theutterance \stepcounter{utterance}  
    & & & \multicolumn{2}{p{0.3\linewidth}}{
        \cellcolor[rgb]{0.9,0.9,0.9}{
            \makecell[{{p{\linewidth}}}]{
                \texttt{\tiny{[GM$|$GM]}}
                \texttt{continue} \\
            }
        }
    }
    & & \\ \\

    \theutterance \stepcounter{utterance}  
    & & \multicolumn{4}{p{0.6\linewidth}}{
        \cellcolor[rgb]{0.9,0.9,0.9}{
            \makecell[{{p{\linewidth}}}]{
                \texttt{\tiny{[GM$\rangle$P2]}}
                \texttt{Dit is mijn BESCHRIJVINGBESCHRIJVING: Mijn raster is een 5x5 vierkant met een dikke rand van "X"{-}symbolen aan alle buitenste randen. In het midden bevindt zich een 3x3 gebied van lege vierkanten (▢), omringd door die rand.} \\
\texttt{Stel nu een vraag om nieuwe aspecten van mijn raster te ontdekken die anders kunnen zijn dan jouw raster. Begin met "VRAAG:" en voeg niets anders toe.} \\
            }
        }
    }
    & \\ \\

    \theutterance \stepcounter{utterance}  
    & & & \multicolumn{4}{p{0.6\linewidth}}{
        \cellcolor[rgb]{0.9,0.9,0.9}{
            \makecell[{{p{\linewidth}}}]{
                \texttt{\tiny{[GM$\langle$P2]}}
                \texttt{VRAAG: Bevatten de hoeken van jouw raster (positie {[}1,1{]}, {[}1,5{]}, {[}5,1{]} en {[}5,5{]}) allemaal een X, of zijn er afwijkingen?} \\
            }
        }
    }
     \\ \\

    \theutterance \stepcounter{utterance}  
    & & & \multicolumn{2}{p{0.3\linewidth}}{
        \cellcolor[rgb]{0.9,0.9,0.9}{
            \makecell[{{p{\linewidth}}}]{
                \texttt{\tiny{[GM$|$GM]}}
                \texttt{continue} \\
            }
        }
    }
    & & \\ \\

    \theutterance \stepcounter{utterance}  
    & & \multicolumn{4}{p{0.6\linewidth}}{
        \cellcolor[rgb]{0.9,0.9,0.9}{
            \makecell[{{p{\linewidth}}}]{
                \texttt{\tiny{[P1$\langle$GM]}}
                \texttt{Dit is mijn BESCHRIJVINGBESCHRIJVING: Een 5x5 raster met een omlijning van X'en, behalve de eerste kolom die volledig uit X'en bestaat. In het midden van het raster (kolommen 2 t/m 4, rijen 2 t/m 4) zijn lege vakjes (▢), met uitzondering van de rechterbovenhoek (rij 1, kolom 4) die een X heeft. De rechteronderhoek (rij 5, kolom 4) is ook een X.} \\
\texttt{VRAAG: Bevatten de hoeken van jouw raster (positie {[}1,1{]}, {[}1,5{]}, {[}5,1{]} en {[}5,5{]}) allemaal een X, of zijn er afwijkingen?Begin je antwoord met "ANTWOORD:" en voeg niets anders toe.} \\
            }
        }
    }
    & \\ \\

    \theutterance \stepcounter{utterance}  
    & \multicolumn{4}{p{0.6\linewidth}}{
        \cellcolor[rgb]{0.9,0.9,0.9}{
            \makecell[{{p{\linewidth}}}]{
                \texttt{\tiny{[P1$\rangle$GM]}}
                \texttt{ANTWOORD: Alle vier de hoeken van mijn raster ({[}1,1{]}, {[}1,5{]}, {[}5,1{]} en {[}5,5{]}) bevatten een X.} \\
            }
        }
    }
    & & \\ \\

    \theutterance \stepcounter{utterance}  
    & & & \multicolumn{2}{p{0.3\linewidth}}{
        \cellcolor[rgb]{0.9,0.9,0.9}{
            \makecell[{{p{\linewidth}}}]{
                \texttt{\tiny{[GM$|$GM]}}
                \texttt{continue} \\
            }
        }
    }
    & & \\ \\

    \theutterance \stepcounter{utterance}  
    & & \multicolumn{4}{p{0.6\linewidth}}{
        \cellcolor[rgb]{0.9,0.9,0.9}{
            \makecell[{{p{\linewidth}}}]{
                \texttt{\tiny{[P1$\langle$GM]}}
                \texttt{Stel nu een vraag om nieuwe aspecten van mijn raster te ontdekken die anders kunnen zijn dan jouw raster. Begin met "VRAAG:" en voeg niets anders toe.} \\
            }
        }
    }
    & \\ \\

    \theutterance \stepcounter{utterance}  
    & \multicolumn{4}{p{0.6\linewidth}}{
        \cellcolor[rgb]{0.9,0.9,0.9}{
            \makecell[{{p{\linewidth}}}]{
                \texttt{\tiny{[P1$\rangle$GM]}}
                \texttt{VRAAG: Bevat jouw raster in de eerste kolom (van boven naar beneden) uitsluitend X'en, of zitten daar ook lege vakjes (▢) tussen?} \\
            }
        }
    }
    & & \\ \\

    \theutterance \stepcounter{utterance}  
    & & & \multicolumn{2}{p{0.3\linewidth}}{
        \cellcolor[rgb]{0.9,0.9,0.9}{
            \makecell[{{p{\linewidth}}}]{
                \texttt{\tiny{[GM$|$GM]}}
                \texttt{continue} \\
            }
        }
    }
    & & \\ \\

    \theutterance \stepcounter{utterance}  
    & & \multicolumn{4}{p{0.6\linewidth}}{
        \cellcolor[rgb]{0.9,0.9,0.9}{
            \makecell[{{p{\linewidth}}}]{
                \texttt{\tiny{[GM$\rangle$P2]}}
                \texttt{ANTWOORD: Alle vier de hoeken van mijn raster ({[}1,1{]}, {[}1,5{]}, {[}5,1{]} en {[}5,5{]}) bevatten een X.} \\
\texttt{VRAAG: Bevat jouw raster in de eerste kolom (van boven naar beneden) uitsluitend X'en, of zitten daar ook lege vakjes (▢) tussen?Begin je antwoord met "ANTWOORD:" en voeg niets anders toe.} \\
            }
        }
    }
    & \\ \\

    \theutterance \stepcounter{utterance}  
    & & & \multicolumn{4}{p{0.6\linewidth}}{
        \cellcolor[rgb]{0.9,0.9,0.9}{
            \makecell[{{p{\linewidth}}}]{
                \texttt{\tiny{[GM$\langle$P2]}}
                \texttt{ANTWOORD: Ja, de eerste kolom van mijn raster bevat uitsluitend X'en (geen lege vakjes).} \\
            }
        }
    }
     \\ \\

    \theutterance \stepcounter{utterance}  
    & & & \multicolumn{2}{p{0.3\linewidth}}{
        \cellcolor[rgb]{0.9,0.9,0.9}{
            \makecell[{{p{\linewidth}}}]{
                \texttt{\tiny{[GM$|$GM]}}
                \texttt{continue} \\
            }
        }
    }
    & & \\ \\

    \theutterance \stepcounter{utterance}  
    & & \multicolumn{4}{p{0.6\linewidth}}{
        \cellcolor[rgb]{0.9,0.9,0.9}{
            \makecell[{{p{\linewidth}}}]{
                \texttt{\tiny{[GM$\rangle$P2]}}
                \texttt{Stel nu een vraag om nieuwe aspecten van mijn raster te ontdekken die anders kunnen zijn dan jouw raster. Begin met "VRAAG:" en voeg niets anders toe.} \\
            }
        }
    }
    & \\ \\

    \theutterance \stepcounter{utterance}  
    & & & \multicolumn{4}{p{0.6\linewidth}}{
        \cellcolor[rgb]{0.9,0.9,0.9}{
            \makecell[{{p{\linewidth}}}]{
                \texttt{\tiny{[GM$\langle$P2]}}
                \texttt{VRAAG: Bevat jouw raster in de tweede rij (van links naar rechts) ook drie opeenvolgende lege vakjes (▢) tussen de X'en aan het begin en einde, of wijkt dit af?} \\
            }
        }
    }
     \\ \\

    \theutterance \stepcounter{utterance}  
    & & & \multicolumn{2}{p{0.3\linewidth}}{
        \cellcolor[rgb]{0.9,0.9,0.9}{
            \makecell[{{p{\linewidth}}}]{
                \texttt{\tiny{[GM$|$GM]}}
                \texttt{continue} \\
            }
        }
    }
    & & \\ \\

    \theutterance \stepcounter{utterance}  
    & & \multicolumn{4}{p{0.6\linewidth}}{
        \cellcolor[rgb]{0.9,0.9,0.9}{
            \makecell[{{p{\linewidth}}}]{
                \texttt{\tiny{[P1$\langle$GM]}}
                \texttt{ANTWOORD: Ja, de eerste kolom van mijn raster bevat uitsluitend X'en (geen lege vakjes).} \\
\texttt{VRAAG: Bevat jouw raster in de tweede rij (van links naar rechts) ook drie opeenvolgende lege vakjes (▢) tussen de X'en aan het begin en einde, of wijkt dit af?Begin je antwoord met "ANTWOORD:" en voeg niets anders toe.} \\
            }
        }
    }
    & \\ \\

    \theutterance \stepcounter{utterance}  
    & \multicolumn{4}{p{0.6\linewidth}}{
        \cellcolor[rgb]{0.9,0.9,0.9}{
            \makecell[{{p{\linewidth}}}]{
                \texttt{\tiny{[P1$\rangle$GM]}}
                \texttt{ANTWOORD: Ja, in de tweede rij van mijn raster zitten drie opeenvolgende lege vakjes (▢) tussen de X'en aan het begin en einde (precies zoals in rij 2, 3 en 4).} \\
            }
        }
    }
    & & \\ \\

    \theutterance \stepcounter{utterance}  
    & & & \multicolumn{2}{p{0.3\linewidth}}{
        \cellcolor[rgb]{0.9,0.9,0.9}{
            \makecell[{{p{\linewidth}}}]{
                \texttt{\tiny{[GM$|$GM]}}
                \texttt{continue} \\
            }
        }
    }
    & & \\ \\

    \theutterance \stepcounter{utterance}  
    & & \multicolumn{4}{p{0.6\linewidth}}{
        \cellcolor[rgb]{0.9,0.9,0.9}{
            \makecell[{{p{\linewidth}}}]{
                \texttt{\tiny{[P1$\langle$GM]}}
                \texttt{Stel nu een vraag om nieuwe aspecten van mijn raster te ontdekken die anders kunnen zijn dan jouw raster. Begin met "VRAAG:" en voeg niets anders toe.} \\
            }
        }
    }
    & \\ \\

    \theutterance \stepcounter{utterance}  
    & \multicolumn{4}{p{0.6\linewidth}}{
        \cellcolor[rgb]{0.9,0.9,0.9}{
            \makecell[{{p{\linewidth}}}]{
                \texttt{\tiny{[P1$\rangle$GM]}}
                \texttt{VRAAG: Bevat jouw raster in de vierde kolom (van links naar rechts) **alleen** de X'en in de hoeken (positie {[}1,4{]} en {[}5,4{]}), of zijn er nog andere X'en in die kolom?} \\
            }
        }
    }
    & & \\ \\

    \theutterance \stepcounter{utterance}  
    & & & \multicolumn{2}{p{0.3\linewidth}}{
        \cellcolor[rgb]{0.9,0.9,0.9}{
            \makecell[{{p{\linewidth}}}]{
                \texttt{\tiny{[GM$|$GM]}}
                \texttt{continue} \\
            }
        }
    }
    & & \\ \\

    \theutterance \stepcounter{utterance}  
    & & \multicolumn{4}{p{0.6\linewidth}}{
        \cellcolor[rgb]{0.9,0.9,0.9}{
            \makecell[{{p{\linewidth}}}]{
                \texttt{\tiny{[GM$\rangle$P2]}}
                \texttt{ANTWOORD: Ja, in de tweede rij van mijn raster zitten drie opeenvolgende lege vakjes (▢) tussen de X'en aan het begin en einde (precies zoals in rij 2, 3 en 4).} \\
\texttt{VRAAG: Bevat jouw raster in de vierde kolom (van links naar rechts) **alleen** de X'en in de hoeken (positie {[}1,4{]} en {[}5,4{]}), of zijn er nog andere X'en in die kolom?Begin je antwoord met "ANTWOORD:" en voeg niets anders toe.} \\
            }
        }
    }
    & \\ \\

    \theutterance \stepcounter{utterance}  
    & & & \multicolumn{4}{p{0.6\linewidth}}{
        \cellcolor[rgb]{0.9,0.9,0.9}{
            \makecell[{{p{\linewidth}}}]{
                \texttt{\tiny{[GM$\langle$P2]}}
                \texttt{ANTWOORD: In de vierde kolom van mijn raster zitten **alleen** X'en in de hoeken ({[}1,4{]} en {[}5,4{]}), en verder geen andere X'en. De posities {[}2,4{]}, {[}3,4{]} en {[}4,4{]} zijn lege vakjes (▢).} \\
            }
        }
    }
     \\ \\

    \theutterance \stepcounter{utterance}  
    & & & \multicolumn{2}{p{0.3\linewidth}}{
        \cellcolor[rgb]{0.9,0.9,0.9}{
            \makecell[{{p{\linewidth}}}]{
                \texttt{\tiny{[GM$|$GM]}}
                \texttt{continue} \\
            }
        }
    }
    & & \\ \\

    \theutterance \stepcounter{utterance}  
    & & \multicolumn{4}{p{0.6\linewidth}}{
        \cellcolor[rgb]{0.9,0.9,0.9}{
            \makecell[{{p{\linewidth}}}]{
                \texttt{\tiny{[GM$\rangle$P2]}}
                \texttt{Stel nu een vraag om nieuwe aspecten van mijn raster te ontdekken die anders kunnen zijn dan jouw raster. Begin met "VRAAG:" en voeg niets anders toe.} \\
            }
        }
    }
    & \\ \\

    \theutterance \stepcounter{utterance}  
    & & & \multicolumn{4}{p{0.6\linewidth}}{
        \cellcolor[rgb]{0.9,0.9,0.9}{
            \makecell[{{p{\linewidth}}}]{
                \texttt{\tiny{[GM$\langle$P2]}}
                \texttt{VRAAG: Bevat jouw raster in de **vijfde kolom** (van links naar rechts) **alleen** de X'en in de hoeken ({[}1,5{]} en {[}5,5{]}), of zijn er ook andere X'en in die kolom (bijv. in rij 2, 3 of 4)?} \\
            }
        }
    }
     \\ \\

    \theutterance \stepcounter{utterance}  
    & & & \multicolumn{2}{p{0.3\linewidth}}{
        \cellcolor[rgb]{0.9,0.9,0.9}{
            \makecell[{{p{\linewidth}}}]{
                \texttt{\tiny{[GM$|$GM]}}
                \texttt{continue} \\
            }
        }
    }
    & & \\ \\

    \theutterance \stepcounter{utterance}  
    & & \multicolumn{4}{p{0.6\linewidth}}{
        \cellcolor[rgb]{0.9,0.9,0.9}{
            \makecell[{{p{\linewidth}}}]{
                \texttt{\tiny{[P1$\langle$GM]}}
                \texttt{ANTWOORD: In de vierde kolom van mijn raster zitten **alleen** X'en in de hoeken ({[}1,4{]} en {[}5,4{]}), en verder geen andere X'en. De posities {[}2,4{]}, {[}3,4{]} en {[}4,4{]} zijn lege vakjes (▢).} \\
\texttt{VRAAG: Bevat jouw raster in de **vijfde kolom** (van links naar rechts) **alleen** de X'en in de hoeken ({[}1,5{]} en {[}5,5{]}), of zijn er ook andere X'en in die kolom (bijv. in rij 2, 3 of 4)?Begin je antwoord met "ANTWOORD:" en voeg niets anders toe.} \\
            }
        }
    }
    & \\ \\

    \theutterance \stepcounter{utterance}  
    & \multicolumn{4}{p{0.6\linewidth}}{
        \cellcolor[rgb]{0.9,0.9,0.9}{
            \makecell[{{p{\linewidth}}}]{
                \texttt{\tiny{[P1$\rangle$GM]}}
                \texttt{ANTWOORD: In de vijfde kolom van mijn raster zitten **alleen** de X'en in de hoeken ({[}1,5{]} en {[}5,5{]}). De posities {[}2,5{]}, {[}3,5{]} en {[}4,5{]} zijn allemaal X'en (dikke rand), dus dat wijkt af van jouw raster.} \\
            }
        }
    }
    & & \\ \\

    \theutterance \stepcounter{utterance}  
    & & & \multicolumn{2}{p{0.3\linewidth}}{
        \cellcolor[rgb]{0.9,0.9,0.9}{
            \makecell[{{p{\linewidth}}}]{
                \texttt{\tiny{[GM$|$GM]}}
                \texttt{continue} \\
            }
        }
    }
    & & \\ \\

    \theutterance \stepcounter{utterance}  
    & & \multicolumn{4}{p{0.6\linewidth}}{
        \cellcolor[rgb]{0.9,0.9,0.9}{
            \makecell[{{p{\linewidth}}}]{
                \texttt{\tiny{[P1$\langle$GM]}}
                \texttt{Stel nu een vraag om nieuwe aspecten van mijn raster te ontdekken die anders kunnen zijn dan jouw raster. Begin met "VRAAG:" en voeg niets anders toe.} \\
            }
        }
    }
    & \\ \\

    \theutterance \stepcounter{utterance}  
    & \multicolumn{4}{p{0.6\linewidth}}{
        \cellcolor[rgb]{0.9,0.9,0.9}{
            \makecell[{{p{\linewidth}}}]{
                \texttt{\tiny{[P1$\rangle$GM]}}
                \texttt{VRAAG: Bevat jouw raster in de **derde rij** (van boven naar beneden) **precies** dezelfde opbouw als de tweede en vierde rij (dus X aan begin en einde, met drie ▢ ertussen), of wijkt dit af?} \\
            }
        }
    }
    & & \\ \\

    \theutterance \stepcounter{utterance}  
    & & & \multicolumn{2}{p{0.3\linewidth}}{
        \cellcolor[rgb]{0.9,0.9,0.9}{
            \makecell[{{p{\linewidth}}}]{
                \texttt{\tiny{[GM$|$GM]}}
                \texttt{continue} \\
            }
        }
    }
    & & \\ \\

    \theutterance \stepcounter{utterance}  
    & & \multicolumn{4}{p{0.6\linewidth}}{
        \cellcolor[rgb]{0.9,0.9,0.9}{
            \makecell[{{p{\linewidth}}}]{
                \texttt{\tiny{[GM$\rangle$P2]}}
                \texttt{ANTWOORD: In de vijfde kolom van mijn raster zitten **alleen** de X'en in de hoeken ({[}1,5{]} en {[}5,5{]}). De posities {[}2,5{]}, {[}3,5{]} en {[}4,5{]} zijn allemaal X'en (dikke rand), dus dat wijkt af van jouw raster.} \\
\texttt{VRAAG: Bevat jouw raster in de **derde rij** (van boven naar beneden) **precies** dezelfde opbouw als de tweede en vierde rij (dus X aan begin en einde, met drie ▢ ertussen), of wijkt dit af?Begin je antwoord met "ANTWOORD:" en voeg niets anders toe.} \\
            }
        }
    }
    & \\ \\

    \theutterance \stepcounter{utterance}  
    & & & \multicolumn{4}{p{0.6\linewidth}}{
        \cellcolor[rgb]{0.9,0.9,0.9}{
            \makecell[{{p{\linewidth}}}]{
                \texttt{\tiny{[GM$\langle$P2]}}
                \texttt{ANTWOORD: Mijn derde rij heeft **precies** dezelfde opbouw als de tweede en vierde rij: een X aan het begin en einde, met drie ▢ ertussen.} \\
            }
        }
    }
     \\ \\

    \theutterance \stepcounter{utterance}  
    & & & \multicolumn{2}{p{0.3\linewidth}}{
        \cellcolor[rgb]{0.9,0.9,0.9}{
            \makecell[{{p{\linewidth}}}]{
                \texttt{\tiny{[GM$|$GM]}}
                \texttt{continue} \\
            }
        }
    }
    & & \\ \\

    \theutterance \stepcounter{utterance}  
    & & \multicolumn{4}{p{0.6\linewidth}}{
        \cellcolor[rgb]{0.9,0.9,0.9}{
            \makecell[{{p{\linewidth}}}]{
                \texttt{\tiny{[GM$\rangle$P2]}}
                \texttt{Neem nu een beslissing. Wat denk je: zijn jouw raster en het andere beschreven raster hetzelfde? Schrijf "BESLISSING: hetzelfde raster" als je denkt dat het hetzelfde raster is of "BESLISSING: verschillende rasters" als je denkt dat het verschillende rasters zijn. Voeg niets anders toe.} \\
            }
        }
    }
    & \\ \\

    \theutterance \stepcounter{utterance}  
    & & & \multicolumn{4}{p{0.6\linewidth}}{
        \cellcolor[rgb]{0.9,0.9,0.9}{
            \makecell[{{p{\linewidth}}}]{
                \texttt{\tiny{[GM$\langle$P2]}}
                \texttt{BESLISSING: verschillende rasters} \\
            }
        }
    }
     \\ \\

    \theutterance \stepcounter{utterance}  
    & & & \multicolumn{2}{p{0.3\linewidth}}{
        \cellcolor[rgb]{0.9,0.9,0.9}{
            \makecell[{{p{\linewidth}}}]{
                \texttt{\tiny{[GM$|$GM]}}
                \texttt{continue} \\
            }
        }
    }
    & & \\ \\

    \theutterance \stepcounter{utterance}  
    & & & \multicolumn{2}{p{0.3\linewidth}}{
        \cellcolor[rgb]{0.9,0.9,0.9}{
            \makecell[{{p{\linewidth}}}]{
                \texttt{\tiny{[GM$|$GM]}}
                \texttt{success} \\
            }
        }
    }
    & & \\ \\

    \theutterance \stepcounter{utterance}  
    & & \multicolumn{4}{p{0.6\linewidth}}{
        \cellcolor[rgb]{0.9,0.9,0.9}{
            \makecell[{{p{\linewidth}}}]{
                \texttt{\tiny{[P1$\langle$GM]}}
                \texttt{ANTWOORD: Mijn derde rij heeft **precies** dezelfde opbouw als de tweede en vierde rij: een X aan het begin en einde, met drie ▢ ertussen.} \\
\texttt{Neem nu een beslissing. Wat denk je: zijn jouw raster en het andere beschreven raster hetzelfde? Schrijf "BESLISSING: hetzelfde raster" als je denkt dat het hetzelfde raster is of "BESLISSING: verschillende rasters" als je denkt dat het verschillende rasters zijn. Voeg niets anders toe.} \\
            }
        }
    }
    & \\ \\

    \theutterance \stepcounter{utterance}  
    & \multicolumn{4}{p{0.6\linewidth}}{
        \cellcolor[rgb]{0.9,0.9,0.9}{
            \makecell[{{p{\linewidth}}}]{
                \texttt{\tiny{[P1$\rangle$GM]}}
                \texttt{BESLISSING: verschillende rasters} \\
            }
        }
    }
    & & \\ \\

    \theutterance \stepcounter{utterance}  
    & & & \multicolumn{2}{p{0.3\linewidth}}{
        \cellcolor[rgb]{0.9,0.9,0.9}{
            \makecell[{{p{\linewidth}}}]{
                \texttt{\tiny{[GM$|$GM]}}
                \texttt{continue} \\
            }
        }
    }
    & & \\ \\

    \theutterance \stepcounter{utterance}  
    & & & \multicolumn{2}{p{0.3\linewidth}}{
        \cellcolor[rgb]{0.9,0.9,0.9}{
            \makecell[{{p{\linewidth}}}]{
                \texttt{\tiny{[GM$|$GM]}}
                \texttt{success} \\
            }
        }
    }
    & & \\ \\

\end{supertabular}
}

\end{document}
